\section{Introduction}
\label{sec:intro}

Driven by the success of large language models (LLMs), vision--language models
(VLMs) that jointly understand images and text have advanced rapidly, led by
systems such as GPT-4V and Gemini. These general-purpose VLMs perform well on
question answering over natural images, but their performance tends to be limited
in domains requiring specialist knowledge, such as medical imaging, because they
lack an adequate grasp of medical terminology and of pathological and
radiological findings. At the same time, retraining large models from scratch is
not realistic in research and clinical environments with constrained GPU
resources. Parameter-efficient fine-tuning (PEFT) methods such as QLoRA
(quantized low-rank adaptation) \citep{dettmers2023qlora} have therefore
attracted attention as an alternative that makes domain-specific fine-tuning
feasible even on consumer-grade GPUs with 16--24\,GB of VRAM.

Applying a PEFT method such as QLoRA in practice, however, requires selecting a
number of hyperparameters --- LoRA rank, target modules, learning rate, and
others --- and the influence of these choices on final performance often remains
unexamined in a systematic way. Automating hyperparameter optimization (HPO) was
identified as an open problem for VLM PEFT in a 2024 NeurIPS survey. Beyond
traditional grid and random search \citep{bergstra2012random} or Bayesian
optimization (Optuna TPE) \citep{akiba2019optuna}, an autonomous search paradigm
has been proposed in which an LLM agent interprets prior experimental results on
its own and proposes the next configuration.

This paper empirically verifies the full process of adapting lightweight VLMs to
the medical imaging VQA (visual question answering) domain on consumer GPUs. We
demonstrate the feasibility of medical VQA domain adaptation for lightweight VLMs
on consumer-class GPUs (executed here on 24\,GB cards; applicable to the 16\,GB
class on the basis of measured VRAM); we systematically analyse how the principal
hyperparameters of QLoRA fine-tuning --- data scale, LoRA rank, and the scope of
target modules --- affect performance; and we verify whether autoresearch-style
autonomous hyperparameter search driven by an LLM agent delivers performance
competitive with Bayesian optimization while providing interpretable search
rationales. These aims are formalized as the three research questions in
Table~\ref{tab:rqs}.

\begin{table}[t]
\centering
\caption{Research questions and null hypotheses.}
\label{tab:rqs}
\small
\begin{tabular}{@{}p{0.05\textwidth}p{0.55\textwidth}p{0.30\textwidth}@{}}
\toprule
\# & Research question & Null hypothesis \\
\midrule
RQ1 & Do lightweight VLMs (2--3B) differ significantly across models in
zero-shot medical VQA performance? & $H_0$: no difference in VQA accuracy across
models \\
RQ2 & Does QLoRA fine-tuning significantly improve medical VQA performance? &
$H_0$: base $=$ fine-tuned performance \\
RQ3 & Does LLM-agent-based autonomous hyperparameter search achieve performance
competitive with Bayesian optimization (Optuna TPE) while providing interpretable
search rationales? & $H_0$: autoresearch $=$ Optuna (TPE) \\
\bottomrule
\end{tabular}
\end{table}

The null hypothesis for RQ3 was set against Optuna (TPE) rather than plain random
search because the justification for adopting LLM-based HPO --- natural-language
explanation of search rationale, use of prior knowledge to understand the
structure of the search space, and traceability of the reasons for configuration
changes --- holds only if such a method can compete with Bayesian optimization,
which is already established as a practical standard, rather than merely being
``better than random.'' Random search is retained solely as a lower-bound
reference. RQ1, RQ2, and RQ3 are verified against measured data in
Sections~4.1, 4.2, and 4.3 respectively, and the discussion cutting across all
three results is consolidated in Section~4.4.

The experimental scope is limited to four lightweight VLMs (Qwen3-VL-2B,
Qwen2.5-VL-3B, SmolVLM2-2.2B, Gemma4-E2B) and three public medical VQA datasets
(PathVQA, SLAKE, VQA-RAD). Models were selected on the criteria that QLoRA
fine-tuning was expected to be feasible within 16\,GB-class VRAM (confirmed
\emph{post hoc} at a measured maximum of 14.4\,GB, Section~\ref{sec:hardware})
and that the license permits free redistribution (Apache 2.0 / MIT). Datasets
were restricted to publicly accessible benchmarks in which clinical question
types are labelled.

\paragraph{Contributions.} This paper makes three contributions.
\begin{itemize}
\itemsep2pt
\item It provides three-stage empirical evidence on medical-domain adaptation of
lightweight VLMs under a single protocol, showing that performance and resource
consumption are not monotonically related and thereby offering a concrete basis
for model selection under resource constraints.
\item It quantifies the model-level heterogeneity of fine-tuning effects and
demonstrates how pooled analysis conceals opposing effects.
\item It reports a negative result for LLM-based autonomous optimization together
with a verification of its failure mechanism. Notably, the variation observed
when repeating an identical configuration (0.056) exceeded the mean difference
between strategies (0.031), constituting a quantitative counterexample to the
practice of comparing methods from single runs.
\end{itemize}
